\documentclass[12pt]{article}
\usepackage[english]{babel}
\usepackage[utf8x]{inputenc}
\usepackage[T1]{fontenc}
\usepackage{scribe}
\usepackage{listings}
\usepackage{quiver}
\usepackage{tikz}


\begin{document}
Songyu Ye

\today
\section{Introduction}
In this note, we will talk a little bit about representations of complex tori 
and $\GL(2,\C)$ as algebraic groups. We will also talk about territories, which
for us will be moduli spaces of curves whose normalizations is $\P^1$.

\section{Representations of Tori}
\subsection{Real tori}
Before people thought about $\Gm(\C)$ as an algebraic group, people thought about $S^1$ as a Lie group.
In this setting, one can ask about the finite dimensional complex representations of $S^1$.

\hfill 

One can also show that all finite dimensional representations of $S^1$ are decomposable, i.e.
can be written as a direct sum of irreducible representations. Thus it is enough to just consider the irreducible representations of $S^1$.

\hfill 

By Schur's Lemma they are all $1$-dimensional and therefore are indexed by characters $\chi: S^1\to \C^*$.
Since $S^1$ is compact, its image in $\C^*$ must also be compact, and we can moreover
it is connected and contains the identity. Therefore, the image of $\chi$ must lie in $S^1$.


\begin{proposition}
All characters of $S^1$ are isomorphic to $\chi_n: S^1\to S^1$ given by $z\mapsto z^n$ for $n\in \Z$.
\end{proposition}

\begin{proof}
    One can prove by making use of the universal covering map $\exp: \R\to S^1$.
    Given a character $\chi: S^1\to S^1$, we can lift it to a map $\tilde \chi: S^1\to \R$.
    Since $\chi$ is a group map, it carries $1$ to $1$, and the fiber over $1$ under $\exp$ is $\Z$.
\end{proof}

In general, a real torus $T$ is a real Lie group whihch is isomorphic to the product of $n$ circles, 
in which case we say that $T$ has rank $n$. Since characters for $S^1\times\cdots\times S^1$ are the same as products of characters for $S^1$, 
all characters of $T$ are indexed by $\Z^n$. Explicitly iff $T$ has rank $n$, then a character $\chi: T\to S^1$ is given by a tuple of integers $(n_1,\ldots, n_k)$, and the character is given by
\begin{align*}
    (z_1,\ldots, z_k)\mapsto z_1^{n_1}\cdots z_k^{n_k}
\end{align*}
From our discussion above, we have the following classification statement for representations of real tori as Lie groups.

\begin{theorem}
    Let $T$ be a real torus of rank $n$. Then every finite dimensional representation $V$ of $T$ is isomorphic 
    to a direct sum of $1$-dimensional irreducible representations with some multiplicities.
    \begin{align*}
        V \cong \bigoplus_{\chi\in \Z^n} W_\chi^{\oplus n_\chi}
    \end{align*} where $W_\chi$ denotes the unique $1$-dimensional irreducible representation for which $T$ acts by $\chi$.
\end{theorem}

In particular, we can decompose $V$ into eigenspaces for the action of $T$. \begin{align*}
    V \cong \bigoplus_{\chi\in\Z^n} V_\chi
\end{align*} where $V_\chi = \set{v\in V \st t\cdot v = \chi(t)v \text{ for all $v\in V$}}$. This is
referred to as the \textbf{weight space decomposition} of $V$ and we refer to the $\chi$ which appear in 
the decomposition as the \textbf{weights} of $V$.


\subsection{Complex tori}
We want an analagous story in algebraic geometry. To get ahold of such a story, 
we need the following framework.
\begin{definition}
    An \textbf{algebraic group} $G$ over $\C$ is an algebraic variety over $\C$ with a group structure
    so that the multipliation map $G\times G\to G$ and inversion map $G\to G$ are morphisms
    of algebraic varieties.
\end{definition}

\begin{definition}
    A \textbf{morphism of algebraic groups} $G\to H$ is a morphism of algebraic varieties that is also a group homomorphism.
\end{definition}

We will consider complex tori $T = \Gm(\C)\times\cdots\times\Gm(\C)$ as algebraic groups, 
in particular the classification statement of rational representations of complex tori 
as algebraic groups. This is an algebraic group because $T$ is the zero locus of 
the polynomial equations \begin{align*}
    T = \Spec \C[x_1^{\pm 1},\ldots, x_n^{\pm 1}] := \Spec\big(\C[x_1,\dots,x_n,y_1,\dots,y_n]/(x_1y_1-1,\dots,x_ny_n-1)\big)
\end{align*} This is a group in the familiar way, and it is clear that the group law is indeed a morphism of algebraic varieties.

\begin{example}
    $\GL(n,\C)$ is a familiar group which is endowed with the structure of an algebraic group.
    $\GL(n,\C)$ is the zero locus of the polynomial equations
    \begin{align*}
        \GL(n,\C) = \Spec \big(\C[x_{ij},\inv{\det}]\big) := \Spec\big(\C[x_{ij},t]/(\det(x_{ij})t-1)\big)
    \end{align*} This variety becomes a group in the familiar way,
    and it is clear that the group law is indeed a morphism of algebraic varieties.
\end{example}

\begin{definition}
    Let $T$ be a complex torus. A \textbf{rational representation} of $T$ is a morphism of algebraic groups
    $T\to \GL(V)$ for some finite dimensional complex vector space $V$.
\end{definition}


\begin{theorem}\label{thm:ratrep}
    Let $T$ be a complex torus.
    Then every finite dimensional rational representation of $T$ 
    decomposes into the direct sum of weight spaces for the action of $T$.
    \begin{align*}
        V \cong \bigoplus_{\chi\in \Z^n} V_\chi
    \end{align*} where $V_\chi = \set{v\in V \st t\cdot v = \chi(t)v \text{ for all $v\in V$}}$
    and $\chi:T\to\GL(1,\C)$ is identified with tuples of integers $(n_1,\ldots, n_k)$ as before.
\end{theorem}
This is an important theorem to the study of representations of algebraic groups.
We will give an proof of this theorem after we introduce the language of Hopf algebras.
\subsection{Hopf algebras}
\begin{definition}
    Let $A$ be a $\C$-algebra. Then we say $A$ is a \textbf{Hopf algebra} if there are maps \begin{align*}
        \text{comultiplication\quad}\Delta: A &\to A\otimes A \\
        \text{counit (augmentation)\quad}\epsilon: A &\to \C \\
        \text{coinverse (antipode)\quad}S: A &\to A
    \end{align*}
    so that the following diagrams commute:
    \begin{center}
        \begin{tikzcd}
            A \arrow[r, "\Delta"] \arrow[d, "\Delta"] & A\otimes A \arrow[d, "\Delta\otimes \id"] \\
            A\otimes A \arrow[r, "\id\otimes \Delta"] & A\otimes A\otimes A
        \end{tikzcd} \\

        \begin{tikzcd}
            A \arrow[r, "\Delta"] \arrow[d, "\id"] & A\otimes A \arrow[d, "\epsilon\otimes \id"] \\
            A \arrow[r, "\cong"] & \C\otimes A
        \end{tikzcd} \\
        
        \begin{tikzcd}
            A \arrow[r, "\Delta"] \arrow[d, "\epsilon"] & A\otimes A \arrow[d, "S \otimes \id"] \\
            \C \arrow[r, ""] & A
        \end{tikzcd}

    \end{center}
\end{definition}
This definition seemingly comes from nowhere, but the following observation will make them more natural
(the names also will make more sense!).

\hfill

Let $G$ be an algebraic group and $\mc O(G)$ its ring of regular functions. Then the multipliation,
inversion, and identity maps \begin{align*}
    \mu: G\times G &\to G \\
    \iota: G &\to G \\
    e: \Spec \C &\to G
\end{align*}
induce maps on the coordinate rings
\begin{align*}
    \Delta: \mc O(G) &\to \mc O(G)\otimes \mc O(G) \\
    \epsilon: \mc O(G) &\to \C \\
    S: \mc O(G) &\to \mc O(G)
\end{align*} where we made the identification $\mc O(G\times G) \cong \mc O(G)\otimes \mc O(G)$.
Because the group axioms hold, these maps satisfy the commutative diagrams above, and therefore
$\mc O(G)$ has the structure of a Hopf algebra.

\begin{remark}
    These maps can be worked out very explicitly. In particular the points of $G$ are in correspondence with the elements of $\Hom_{\kAlg}(\cO(G),\C)$.
    The correspondence is very explicitly given by $g\mapsto \ev_g$ where $\ev_g: \cO(G)\to \C$ is the evaluation map. 
    The key idea is that for an arbitrary algebraic group $G$ and two points $x,y$ of $G$,
     $f,g: \cO(G)\to \C$ the corresponding morphisms of $\C$-algebras, then the composition $(f\otimes g)\circ \Delta$ is again a map $\cO(G)\to \cO(G)\otimes \cO(G)\to \C$.
    The condition that we ask from comultiplication is precisely that the composition $(f\otimes g)\circ \Delta$
    corresponds to the product $xy\in G$.
\end{remark}

\begin{example}
Recall that
    \begin{align*}
        \cO(\Ga(\C)) &= \C[x]
    \end{align*} where $\Ga(\C)$ is the additive group of $\C$. 
    Let $f,g\in \Hom_{\kAlg}(\cO(G),\C)$ with $f(x) = a$ and $g(x) = b$. 
    We want to find a map $\Delta: \cO(\Ga(\C))\to \cO(\Ga(\C))\otimes \cO(\Ga(\C))$ so that
    \begin{align*}
        ((f\otimes g)\circ \Delta)(X) = (a+b)
    \end{align*}
    We write down the map $\Delta$ explicitly as
    \begin{align*}
        \Delta(X) = X\otimes 1 + 1\otimes X
    \end{align*} and notice that it does the job. We see that $\Delta$ then must be the 
    comultiplication map for $\Ga(\C)$ because such map is unique, given the prescribed group law on $\Ga(\C)$.
\end{example}

\begin{example}
    By the same token, we can work out the Hopf algebra structure for $\Gm(\C)$ to be
    \begin{align*}
        \Delta(x) &= x\otimes x \\
        \epsilon(x) &= 1 \\
        S(x) &= \inv{x}
    \end{align*}
\end{example}

\begin{example}
    Consider the example of $\GL(2,\C)$ as an algebraic group. The Hopf algebra structure is given by
    \begin{align*}
        \Delta(x_{ij}) &= \sum_{k=1}^2 x_{ik}\otimes x_{kj} \\
        \epsilon(x_{ij}) &= \delta_{ij} \\
        S(x_{ij}) &= M_{ij}
    \end{align*} where \begin{align*}
        M = \begin{bmatrix}
            x_{11} & x_{12} \\
            x_{21} & x_{22}
        \end{bmatrix}^{-1} = \frac{1}{\det(M)}\begin{bmatrix}
            x_{22} & -x_{12} \\
            -x_{21} & x_{11}
        \end{bmatrix}
    \end{align*}
\end{example}

\hfill

Now we want to translate the representation theory of algebraic groups $G$ into the language of comodules over Hopf algebras.
\begin{theorem}
    Let $G$ be an algebraic group. Then rational representations $V$ of $G$ 
    correspond to linear maps $\rho: V\to V\otimes \mc O(G)$ satisfying the following diagrams:
    \begin{center}
        \begin{tikzcd}
            V \arrow[r, "\rho"] \arrow[d, "\rho"] & V\otimes \mc O(G) \arrow[d, "\rho\otimes \id"] \\
            V\otimes \mc O(G) \arrow[r, "\id\otimes \Delta"] & V\otimes \mc O(G)\otimes \mc O(G)
        \end{tikzcd} \\

        \begin{tikzcd}
            V \arrow[r, "\rho"] \arrow[d, "\id"] & V\otimes \mc O(G) \arrow[d, "\id\otimes \epsilon"] \\
            V \arrow[r, "\cong"] & V\otimes \C
        \end{tikzcd} \\
    \end{center}
\end{theorem}
\begin{proof}
    [Sketch of proof] The first diagram says that the action of $G$ on $V$ is associative 
    and the second diagram says that $e\in G$ acts by the identity transformation on $V$. 
    These are precisely the conditions that say that $V$ is a representation of $G$.
\end{proof}
We refer the reader to \cite{waterhouse} for a more detailed discussion of this theorem.

\begin{definition}
    We call such a linear map $\rho$ a \textbf{comodule structure} on $V$.
\end{definition}

\begin{example}
    Consider the action of $\Gm(\C)$ on $\C^2$ given by
    \begin{align*}
        t\cdot (a,b) = (ta, t^{-1}b)
    \end{align*} This is a rational representation of $\Gm(\C)$ which we can write as 
    \begin{align*}
        \tau:\Gm(\C) &\to \GL(2,\C) \\
        t &\mapsto \begin{bmatrix}
            t & 0 \\
            0 & \inv{t}
        \end{bmatrix}
    \end{align*} This induces a comodule structure on $\C^2$ given by the map $\rho: \C^2\to \C^2\otimes \mc O(\Gm(\C))$ given by
    \begin{align*}
        \rho(a) &= a\otimes x \\
        \rho(b) &= b\otimes \inv{x}
    \end{align*} where $x$ is the coordinate function on $\Gm(\C)$.
\end{example}

\subsection{Weight space decomposition}
We are now ready to give a proof of Theorem \ref{thm:ratrep} using the language of Hopf algebras.
\begin{proof}
    [Proof of \ref{thm:ratrep}]
\end{proof}

\section{Representations of $\GL(2,\C)$}
\section{Territories}

\begin{thebibliography}{9}

\bibitem{waterhouse}
Waterhouse, W. C. (1996). \textit{Introduction to Affine Group Schemes}. Springer.

\bibitem{milne}
Milne, J. S. (2017). \textit{Algebraic Groups: The Theory of Group Schemes of Finite Type over a Field}. Cambridge University Press.

\end{thebibliography}
\end{document}